%%%%%%%%%%%%%%%%%%%%%%%%%%%%%%%%%%%%%%
% Main Preamble
%%%%%%%%%%%%%%%%%%%%%%%%%%%%%%%%%%%%%%
\documentclass[10pt]{article}
\usepackage{fullpage}
\usepackage{amsfonts}
\usepackage{amsmath}
\usepackage{amsthm}
\usepackage{graphicx}
\usepackage{color}
\usepackage{amssymb}
\usepackage{empheq}
\usepackage{mathrsfs}
\usepackage{enumerate}
\usepackage{tikz}
\usepackage{pgflibraryarrows}
\usepackage{pgflibrarysnakes}
\usepackage{upgreek}
\usepackage{tipa}
\usepackage{multicol}
\usepackage{verbatim}
\usepackage{floatrow}
\usepackage{gensymb}
\usepackage{caption}

\usepackage[T1]{fontenc}
\usepackage[font=small,labelfont=bf,tableposition=top]{caption}

\DeclareCaptionLabelFormat{andtable}{#1~#2  \&  \tablename~\thetable}

%%%%%%%%%%%%%%%%%%%%%%%%%%%%%%%%%%%%%%
% Header/footer
%%%%%%%%%%%%%%%%%%%%%%%%%%%%%%%%%%%%%%
\usepackage{fancyhdr}
\setlength{\headheight}{15.2pt}
\pagestyle{fancy}
\setlength\headsep{30pt}
\lhead{P2.WS3}
\rhead{Kutztown University -- MATH 210}
\fancyfoot{}

%%%%%%%%%%%%%%%%%%%%%%%%%%%%%%%%%%%%%%
% Document
%%%%%%%%%%%%%%%%%%%%%%%%%%%%%%%%%%%%%%
\begin{document}

\begin{center}
\Large{\textsc{Phase 2 (part 3): (Non)Linear Curve-Fitting}}\\
\end{center}

\begin{enumerate}[{$1.]$}]

\item A lot of real-world data is exponential in nature and so fitting a function of the form $f(x) =  ae^{bx}$ is necessary.  Show that for any set of data, fitting this curve can be transformed into a linear fitting problem. \vfill


\item A lot of real-world data is logistic in nature and so fitting a function of the form $f(x) = \displaystyle \frac{1}{1 + ae^{-bx}}$ is necessary.  Here, $a$ and $b$ are both strictly positive.  Show that for any set of data, fitting this curve can be transformed into a linear fitting problem. \vfill



\pagebreak

\item Similar to the previous example, we may want to fit a ``Witch of Agnesi" function to specific data.  This function has the form $\displaystyle f(x) = \frac{a}{bx^2 + 1}$, and looks similar to a Normal curve.  Here, $a>0$ is a parameter that controls the height of the centered peak, and $b>0$ controls the width.  Show that a function of this form can be fit to data by transforming it into a linear problem. \vfill



\item Data may following a periodic or sinusoidal pattern.  Therefore, fitting sinusoidal data to a function such as 
\[ f(x) = c_0 + c_1\sin(x) + c_2\cos(x) + c_3\sin(2x) + c_4\cos(2x), \]
may be beneficial.  Show that fitting this function to data is also a linear least-squares problem. \vfill





\end{enumerate}

\end{document}